\documentclass[twocolumn]{iopjournal}

\usepackage{amsmath,amssymb,amsfonts}
\usepackage{array}

\fancyhfoffset[L]{0mm}
\setlength{\marginparwidth}{0pt}
\setlength{\marginparsep}{0pt}

% ==========================================
% ZENODO PREPRINT HEADERS (ACTIVE)
% ==========================================
\renewcommand{\articletype}[1]{%
  {\vspace*{-8mm}\noindent \Large \sf Preprint}\par
  \vspace*{6mm}\noindent{\scriptsize \sf{\bfseries \MakeUppercase{#1}}}}

\fancyhead[L]{}
\fancyhead[R]{}

\begin{document}

\articletype{Supplementary Material}

\title{The Spatial Geometry of the Lorentz Tetrad}

\author{Stephen Nelson}

\affil{Field Dynamics Research}

\ead{stephen@fielddynamics.org}

\begin{abstract}
The legs of a Lorentz tetrad have equal magnitude and are pairwise
orthogonal, so the magnitude pattern $|\eta_{ab}| = \delta_{ab}$ of
the Minkowski metric is invariant under every permutation of the
four legs; only the sign of $\eta_{00}$ distinguishes the time leg.
We realise the four legs as vectors in the three-dimensional
spatial sector and ask which configurations are possible. Three
premises, each read from $\eta$, suffice: the Gram matrix of the
realisation inherits the permutation symmetry of $|\eta|$; the legs
are distinct; and the realisation lives in $\mathbb{R}^3$. These
force the Gram matrix $G = c^2 I - (c^2/3)(J - I)$, and hence the
regular tetrahedron, unique up to $O(3)$. Equal norms and closure
$\sum_a v_a = 0$ are derived, not assumed. Three is the only
dimension in which such a realisation both exists and is forced. A
second route, which grants only the symmetry of $\eta$ itself, gives
the same tetrahedron once closure is assumed. Each premise is shown
to be necessary. Parity completion gives the cube and the
stellated octahedron.
\end{abstract}

% ----------------------------------------------------------------------

\section*{Setting}

A Lorentz tetrad has four legs $e_a$, $a = 0,1,2,3$, with
$\eta(e_a, e_b) = \eta_{ab} = \mathrm{diag}(-1,+1,+1,+1)$. Every leg
has unit magnitude and every pair is orthogonal, so
\begin{equation}
  |\eta_{ab}| = \delta_{ab},
  \label{eq:abseta}
\end{equation}
which is invariant under all $24$ permutations of the legs. The
metric $\eta$ itself is invariant only under the $6$ permutations
that fix the time leg; the sign of $\eta_{00}$ is the only datum
that distinguishes it.

A \emph{spatial realisation} is a $4$-tuple
$V = (v_0, v_1, v_2, v_3)$ of vectors $v_a \in \mathbb{R}^D$, one
per leg, with Gram matrix $G_{ab} = v_a \cdot v_b$. For
$\sigma \in S_4$ let $P_\sigma$ be its permutation matrix. We write
$I$ for the $4 \times 4$ identity, $J$ for the all-ones matrix and
$\mathbf{1} = (1,1,1,1)^T$.

% ----------------------------------------------------------------------

\section*{Premises}

\paragraph{(S) Equivalence of the legs.}
$P_\sigma G P_\sigma^T = G$ for every $\sigma \in S_4$. A Gram
matrix of real vectors is positive semidefinite, so its diagonal
cannot carry the sign $\eta_{00} = -1$; what a realisation can
inherit from $\eta$ is its magnitude pattern~\eqref{eq:abseta}, and
with it the permutation symmetry of the legs.

\paragraph{(N) Distinct legs.} The $v_a$ are not all equal, as the
legs of a tetrad are linearly independent ($\det e \neq 0$).

\paragraph{(D) Dimension.} $D = 3$: the realisation lives in the
spatial sector, the three positive entries of $\eta$.

\medskip
\noindent
Neither equal norms nor closure is assumed.

% ----------------------------------------------------------------------

\section*{Two Lemmas}

\paragraph{Lemma 1 (commutant).} \textit{$G$ satisfies (S) if and
only if $G = \alpha I + \beta J$ for some $\alpha, \beta \in
\mathbb{R}$.}

\paragraph{Proof.} (S) states $G_{\sigma(a)\sigma(b)} = G_{ab}$ for
all $\sigma$, $a$, $b$, so $G$ is constant on the orbits of $S_4$ on
ordered pairs. There are two: the four diagonal pairs, and the
twelve off-diagonal pairs, one orbit because $S_4$ is
$2$-transitive. Writing the two values as $\alpha + \beta$ and
$\beta$ gives $G = \alpha I + \beta J$; conversely $I$ and $J$
commute with every $P_\sigma$. \hfill $\square$

\paragraph{Lemma 2 (rigidity).} \textit{If $U, W \in
\mathbb{R}^{D \times 4}$ have $U^T U = W^T W$ of rank $D$, then
$W = QU$ for a unique $Q \in O(D)$.}

\paragraph{Proof.} For every $\lambda \in \mathbb{R}^4$,
$|U\lambda|^2 = \lambda^T U^T U \lambda = |W\lambda|^2$. Hence
$U\lambda = U\lambda'$ implies $W\lambda = W\lambda'$, and
$Q : U\lambda \mapsto W\lambda$ is a well-defined linear isometry
from the column space of $U$ onto that of $W$. Both are
$\mathbb{R}^D$ because the rank is $D$, so $Q \in O(D)$, and it is
fixed by its values on the columns. \hfill $\square$

% ----------------------------------------------------------------------

\section*{The Forced Gram Matrix}

\paragraph{Theorem.} \textit{Under (S), (N) and (D), the legs have a
common norm $c > 0$, they close, $\sum_a v_a = 0$, and}
\begin{equation}
  \boxed{G = c^2 I - \frac{c^2}{3}\,(J - I),}
  \label{eq:forced}
\end{equation}
\textit{that is $G_{aa} = c^2$ and $G_{ab} = -c^2/3$ for
$a \neq b$.}

\paragraph{Proof.} By Lemma 1, $G = \alpha I + \beta J$. Since
$J\mathbf{1} = 4\mathbf{1}$ and $J$ vanishes on
$\mathbf{1}^\perp$, $G$ has eigenvalue $\alpha + 4\beta$ on
$\mathbf{1}$ and $\alpha$, three times, on $\mathbf{1}^\perp$. Every
pair of legs is at the same distance,
$|v_a - v_b|^2 = G_{aa} + G_{bb} - 2G_{ab} = 2\alpha$, so by (N)
$\alpha > 0$ and $\operatorname{rank} G \ge 3$. By (D),
$\operatorname{rank} G \le 3$. Hence $\alpha + 4\beta = 0$, so
$G\mathbf{1} = 0$ and
$|\sum_a v_a|^2 = \mathbf{1}^T G \mathbf{1} = 0$: the legs close.
Their common norm is $c^2 = G_{aa} = \alpha + \beta = 3\alpha/4$,
and $\beta = -\alpha/4 = -c^2/3$, which is
Eq.~\eqref{eq:forced}. \hfill $\square$

\paragraph{Corollary 1 (the regular tetrahedron).} \textit{The four
points $v_a$ are pairwise equidistant, with edge $c\sqrt{8/3}$ and
centroid at the origin: they are the vertices of a regular
tetrahedron of circumradius $c$. With $c = \sqrt3$ the columns of}
\begin{equation}
  M = \begin{pmatrix}
    +1 & +1 & -1 & -1 \\
    +1 & -1 & +1 & -1 \\
    +1 & -1 & -1 & +1
  \end{pmatrix}
  \label{eq:M}
\end{equation}
\textit{realise Eq.~\eqref{eq:forced}, and every realisation is
$V = (c/\sqrt3)\,QM$ for a unique $Q \in O(3)$.}

\paragraph{Proof.} The distances are $2\alpha = 8c^2/3$ and closure
puts the centroid at the origin. Each column of $M$ has three
entries $\pm1$, and each pair of columns agrees in exactly one
place, so $M^TM = 3I - (J - I)$. Both $V$ and $(c/\sqrt3)M$ have
Gram matrix~\eqref{eq:forced}, of rank $3$, and Lemma 2 gives $Q$.
\hfill $\square$

\paragraph{Corollary 2 (why three dimensions).} \textit{For $K$ legs
satisfying (S) and (N) in $\mathbb{R}^D$, a realisation exists if
and only if $D \ge K - 1$, and it is forced to the closed regular
simplex if and only if $D = K - 1$. For the tetrad, $K = 4$ and
$D = 3$.}

\paragraph{Proof.} The argument above with $K$ in place of $4$ gives
eigenvalues $\alpha > 0$ ($K - 1$ times) and $\alpha + K\beta$, so
$\operatorname{rank} G \ge K - 1$ and no realisation exists for
$D < K - 1$. The regular simplex lies in $\mathbb{R}^{K-1} \subseteq
\mathbb{R}^D$, so one exists for $D \ge K - 1$. If $D = K - 1$, then
$\alpha + K\beta = 0$ and the simplex is forced. If $D \ge K$, any
$\alpha + K\beta > 0$ is also realisable, for example $\beta = 0$:
$K$ orthogonal legs of equal length. \hfill $\square$

\paragraph{Corollary 3 (parity completion).} \textit{Spatial parity
$v_a \to -v_a$ preserves $G$, so $-V$ is a second realisation. $V$
and $-V$ share no point, and for $V = M$ their eight points are the
vertices of the cube $\{-1,+1\}^3$; the compound of the two
tetrahedra is Kepler's stella
octangula~\cite{Kepler1619,Cromwell1997}, the stellated
octahedron.}

\paragraph{Proof.} $(-V)^T(-V) = V^TV$. A shared point
$v_a = -v_b$ would need $G_{ab} = -c^2$, but $G_{ab} \in
\{c^2, -c^2/3\}$. The columns of $M$ and $-M$ are eight distinct
sign triples, which exhaust $\{-1,+1\}^3$. \hfill $\square$

% ----------------------------------------------------------------------

\section*{A Second Route}

A reader may grant only the symmetry of $\eta$ itself, the
permutations of the three spatial legs, together with equal norms.
Closure must then be assumed, and the result is the same.

\paragraph{Proposition.} \textit{Let $G$ be invariant under the
permutations of legs $1, 2, 3$, with $G_{aa} = c^2 > 0$ for all
$a$, and let $\sum_a v_a = 0$. Then $G$ is
Eq.~\eqref{eq:forced}.}

\paragraph{Proof.} The invariance gives $G_{0i} = q$ and
$G_{ij} = t$ for $i \neq j$ in $\{1,2,3\}$. Closure gives
$G\mathbf{1} = 0$. Row $0$ reads $c^2 + 3q = 0$, and row $i$ reads
$q + c^2 + 2t = 0$, so $q = t = -c^2/3$. \hfill $\square$

\medskip
\noindent
The time leg, which $\eta$ singles out by its sign, ends up
geometrically equivalent to the spatial legs: the realisation
recovers the full permutation symmetry of $|\eta|$.

% ----------------------------------------------------------------------

\section*{Necessity of the Premises}

Removing any one premise admits configurations that satisfy the
others but are not the regular tetrahedron:

\vspace{6pt}
\noindent
\begin{tabular*}{\columnwidth}{@{\extracolsep{\fill}}l>{\raggedright\arraybackslash}p{0.74\columnwidth}@{}}
\hline
Removed & Admitted configurations \\
\hline
(S) & $(v, w, -v, -w)$ with $|v| = |w|$, $v \perp w$: distinct legs
  in $\mathbb{R}^3$, equal norms and even closed, but a square \\
(N) & $v_a = v$ for all $a$: $G = c^2 J$ \\
(D) & In $\mathbb{R}^4$: $v_a = c\,\hat e_a$, $G = c^2 I$, the
  pattern $|\eta|$ itself, not closed. In $\mathbb{R}^2$: no
  realisation \\
Closure & (second route) $v_i = c\,\hat e_i$,
  $v_0 = -c(1,1,1)/\sqrt3$: $G_{0i} = -c^2/\sqrt3$, $G_{ij} = 0$ \\
\hline
\end{tabular*}
\vspace{6pt}

In the second route, invariance under the spatial legs with equal
norms leaves, for distinct legs in $\mathbb{R}^3$, the one-parameter
family $c^2(c^2 + 2t) = 3q^2$; closure selects $q = t = -c^2/3$ from
it.

% ----------------------------------------------------------------------

\section*{Summary}

Everything is read from the Minkowski metric. Its magnitude pattern
$|\eta_{ab}| = \delta_{ab}$ makes the legs equivalent, and its three
positive entries fix the dimension of the spatial sector. In that
sector, four equivalent and distinct legs have exactly one
configuration up to $O(3)$: the regular tetrahedron, with common
norm and closure as consequences. Three is the only dimension in
which this configuration both exists and is forced. Parity, which
the Gram matrix cannot see, supplies the second tetrahedron, and
the two together form the stellated octahedron.

% ----------------------------------------------------------------------

\begin{thebibliography}{9}

\bibitem{Kepler1619}
J.~Kepler,
\textit{Harmonices Mundi}
(Linz, 1619).

\bibitem{Cromwell1997}
P.~R.~Cromwell,
\textit{Polyhedra}
(Cambridge University Press, Cambridge, 1997).

\end{thebibliography}

\end{document}
